\documentclass[]{article}

\usepackage{amsmath}
\usepackage{multirow}
\usepackage{natbib}
\usepackage{graphicx} 


\begin{document}

Distinguishing the determinants of the investment rate from the factors that moderate the speed of economic convergence is a central challenge in growth economics. This study develops a dynamic panel framework to isolate these distinct mechanisms using a synthetic dataset for 50 countries (1990–2019) generated from a structural growth model. We first employ a System Generalized Method of Moments (GMM) estimator to quantify the endogenous feedback between income levels and the investment rate. Subsequently, we model capital accumulation as a function of the distance to a calculated steady-state, using instrumented interaction terms to test whether trade openness and government expenditure share accelerate or dampen the speed of convergence. Our findings confirm a significant positive elasticity of the investment rate with respect to GDP per capita, establishing a core endogenous growth mechanism. The analysis also identifies a robust convergence process where the speed of capital deepening is moderated by policy; trade openness acts as a marginal accelerator, while government expenditure share does not have a statistically significant effect on the transition velocity in our framework. These results demonstrate a clear distinction between the income-driven determinants of investment levels and the policy-related factors that influence the dynamics of capital accumulation.

\end{document}
                